\documentclass[11pt]{article}
\usepackage[utf8]{inputenc}
\usepackage{amsmath,amssymb,amsthm}
\usepackage{hyperref}
\usepackage{geometry}
\geometry{margin=1in}

\title{Refutation of TLMC Conjecture 00000004225\\
\large Higher Dirichlet--Neumann zero-multiplicity identity}
\author{orionsheep (Lean 4 + Mathlib formalization)}
\date{2026-10-03}

\begin{document}
\maketitle

\section*{The conjecture}
\noindent\textbf{Definition (as stated).} The higher Dirichlet--Neumann operator
is the boundary response operator on $k$-chains.

\noindent\textbf{Conjecture.} The multiplicity of its zero eigenvalue equals
exactly the count of $k$-spanning trees of the boundary, and the multiplicity
is one if and only if the boundary is a sphere triangulation.

\section*{Standard reading and why it fails}
On a boundary complex $B=\partial M$, the combinatorial higher
Dirichlet--Neumann (boundary-response) operator on $k$-chains is
\[
\Delta_k \;=\; \partial_{k+1}\partial_{k+1}^{\!\top}
             \;+\; \partial_k^{\!\top}\partial_k ,
\]
the Hodge Laplacian. Its zero eigenspace is the space of harmonic $k$-chains,
so
\[
\operatorname{mult}_k(B)=\dim\ker\Delta_k(B)=\beta_k(B)
   = f_k-\operatorname{rank}\partial_k-\operatorname{rank}\partial_{k+1}.
\]
The simplicial spanning-tree count $\tau_k(B)$ is a \emph{determinantal}
quantity (Kalai's matrix-tree theorem: a weighted product of the
\emph{nonzero} Laplacian eigenvalues), not a kernel dimension. Hence
$\operatorname{mult}_k=\tau_k$ fails generically, and the ``multiplicity one
iff sphere'' clause fails in both directions.

\section*{Counterexample A: $\partial\Delta^3$ (a sphere), $k=1$}
The boundary of the tetrahedron has $(f_0,f_1,f_2)=(4,6,4)$, $\chi=2$: a
triangulation of $S^2$. Direct computation gives
\[
\Delta_1 = \partial_1^{\top}\partial_1+\partial_2\partial_2^{\top}=4I_6,
\qquad \det\Delta_1=4096,
\]
so $\operatorname{mult}_1=0$. But $\tau_1=\tau(K_4)=16$ (Cayley). Thus
$\operatorname{mult}_1\neq\tau_1$ (first clause fails), and a sphere boundary
has $\operatorname{mult}_1\neq 1$ (the ``if'' direction of clause 2 fails).

\section*{Counterexample B: the 7-vertex torus $T^2$, $k=2$}
The minimal triangulation of the torus (the boundary of a solid torus) has
$(f_0,f_1,f_2)=(7,21,14)$, $\chi=0\neq2$: not a sphere. Since there are no
$3$-faces, $\Delta_2=\partial_2^{\top}\partial_2$ on $C_2\cong\mathbb{Q}^{14}$.
An explicit harmonic $2$-cycle
$c=(-1,\dots,-1,+1,\dots,+1)$ lies in the kernel, and a dual-tree elimination
on the 21 edge equations shows the kernel is exactly $\operatorname{span}\{c\}$;
hence $\operatorname{mult}_2=\beta_2(T^2)=1$. A non-sphere boundary therefore
attains multiplicity $1$ (the ``only if'' direction fails). Moreover
$\tau_2=0$ (a $2$-tree would require $\operatorname{rank}\partial_2=f_2=14$,
but $\operatorname{rank}\partial_2=13$), so clause 1 fails again:
$1\neq0$.

\section*{Formalization}
All objects are defined concretely in \texttt{Main.lean} (namespace
\texttt{Submission00000004225}): \texttt{dnOperator},
\texttt{zeroMult}, \texttt{FinSurface} with \texttt{IsTwoSphere},
the incidence matrices \texttt{tetB1}, \texttt{tetB2}, \texttt{torusB2},
and the tree count \texttt{tetTreeCount}. The headline theorem
\texttt{disproof\_00000004225} states
\begin{align*}
&\operatorname{zeroMult}(\texttt{tetDN1})\neq\texttt{tetTreeCount}\\
&{}\wedge\ \bigl(\texttt{tetBoundary.IsTwoSphere}
   \wedge\operatorname{zeroMult}(\texttt{tetDN1})\neq1\bigr)\\
&{}\wedge\ \bigl(\neg\texttt{torusBoundary.IsTwoSphere}
   \wedge\operatorname{zeroMult}(\texttt{torusDN2})=1\bigr),
\end{align*}
refuting both clauses and both directions. Every proof is axiom-audited in
\texttt{Check.lean}; only the ambient axioms \texttt{propext},
\texttt{Classical.choice}, \texttt{Quot.sound} appear. Numerical
verification is in \texttt{reproduce.py} (sympy, exact arithmetic).

\end{document}
